\documentclass[../../../include/open-logic-section]{subfiles}

\begin{document}

\olfileid{sth}{z}{power}
\olsection{مجموعه‌های توانی}

با اصل موضوع دیگری ادامه می‌دهیم:

\begin{axiom}[مجموعه‌های توانی]
برای هر مجموعهٔ~$A$، مجموعهٔ~$\Pow{A} = \Setabs{x}{x \subseteq A}$ وجود دارد.
\[
	\forall A \exists P \forall x(x \in P \liff (\forall z \in x)z \in A)
\]
\end{axiom}

توجیه ما برای این اصل کاملاً سرراست است. فرض کنید~$A$ در مرحلهٔ~$S$
تشکیل شده باشد. آنگاه همهٔ اعضای~$A$ پیش از~$S$ در دسترس بوده‌اند (بنا
بر~\stagesacc). پس با استدلالی مانند توجیه جداسازی، هر زیرمجموعهٔ~$A$ تا
مرحلهٔ~$S$ تشکیل شده است. بنابراین، همهٔ آن‌ها در دسترس‌اند تا در هر
مرحلهٔ پس از~$S$ به صورت مجموعه‌ای واحد تشکیل شوند. و می‌دانیم چنین
مرحله‌ای وجود دارد، زیرا~$S$ مرحلهٔ آخر نیست (بنا بر~\stagessucc). پس
$\Pow{A}$ وجود دارد.

در ادامه پیامد خوبی از اصل مجموعه‌های توانی می‌آید:

\begin{prop}\label{thm:Products}
برای هر دو مجموعهٔ~$A,B$، حاصل‌ضرب دکارتی آن‌ها، $A \times B$، وجود دارد.
\end{prop}

\begin{proof}
مجموعهٔ~$\Pow{\Pow{A \cup B}}$ بنا بر اصل مجموعه‌های توانی و
\olref[sth][z][pairs]{prop:pairsconsequences} وجود دارد. پس با جداسازی،
مجموعهٔ زیر وجود دارد:
\[
	C = \Setabs{z \in \Pow{\Pow{A \cup B}}}{(\exists x \in A)(\exists y \in B) z = \tuple{x, y}}.
\]
اکنون برای هر~$x \in A$ و~$y \in B$، مجموعهٔ~$\tuple{x,y}$ بنا
بر~\olref[sth][z][pairs]{prop:pairsconsequences} وجود دارد. افزون بر این،
چون~$x,y \in A \cup B$، داریم~$\{x\},\{x,y\} \in \Pow{A \cup B}$ و
$\tuple{x,y} \in \Pow{\Pow{A \cup B}}$. پس~$A \times B = C$.
\end{proof}

در این اثبات، اصل مجموعه‌های توانی با جداسازی تعامل دارد، و این
شگفت‌آور نیست. بدون جداسازی، اصل مجموعه‌های توانی اصل چندان
\emph{نیرومندی} نبود. زیرا جداسازی به ما می‌گوید کدام زیرمجموعه‌های یک
مجموعه وجود دارند و در نتیجه تعیین می‌کند هر مجموعهٔ توانی تا چه اندازه
«فربه» است.

\begin{prob}
نشان دهید برای هر دو مجموعهٔ~$A,B$: (i) مجموعهٔ همهٔ رابطه‌هایی که دامنهٔ
آن‌ها~$A$ و بردشان~$B$ است وجود دارد؛ و (ii) مجموعهٔ همهٔ تابع‌ها از~$A$
به~$B$ وجود دارد.
\end{prob}

\begin{prob}
فرض کنید~$A$ مجموعه و~$\sim$ رابطه‌ای هم‌ارزی روی~$A$ باشد. اثبات کنید
مجموعهٔ کلاس‌های هم‌ارزی تحت~$\sim$ روی~$A$، یعنی
$\equivclass{A}{\sim}$، وجود دارد.
\end{prob}

\end{document}
